Este capítulo retoma o percurso do trabalho em três movimentos. O primeiro verifica, objetivo por objetivo, o que foi efetivamente entregue. O segundo reúne os resultados numéricos que sustentam as escolhas de projeto. O terceiro delimita o que ficou fora do alcance e indica o que a continuidade do trabalho exigiria.

\section{Alcance dos objetivos}\label{sec:7.1}

O primeiro objetivo específico consistia em levantar, na literatura, os parâmetros arquiteturais de uma memória cache e justificar as escolhas adotadas. A seção~\ref{sec:2.5.1} respondeu à parte descritiva, e a seção~\ref{sec:4.2} fixou a configuração do projeto em 8 KiB por cache, linha de 16 bytes, quatro vias e política LRU. A justificativa, no entanto, apoiava-se em duas fontes de terceiros: a arquitetura de referência CVA6 fornecia os valores, e os dados de síntese do BRISC-V indicavam o custo em área de cada parâmetro. Faltava o comportamento do próprio subsistema sob carga. Essa lacuna foi preenchida no \refCap{6_Verificacao}, onde a varredura mostra que 8 KiB e quatro vias se situam exatamente no ponto em que o ganho marginal se anula. Por conseguinte, o objetivo está alcançado com uma justificativa mais forte do que a inicialmente prevista.

O segundo objetivo pedia a especificação da máquina de estados que governa acertos e faltas. O \refCap{4_Especificacao} a fixou em oito estados, com as condições de transição e os protocolos de troca com o pipeline e com a memória descritos em termos de sinais e de temporização, sem recurso a construtos da linguagem. A separação entre o caminho de acesso e o caminho de consistência, com cinco e três estados respectivamente, decorreu dessa especificação. O objetivo está alcançado.

O terceiro objetivo tratava da implementação das duas caches como módulos distintos, com a observação de que a cache de instrução dispensa o bit de sujeira. A seção~\ref{sec:5.4} descreve a solução adotada: uma única classe, instanciada duas vezes com um parâmetro de elaboração que suprime o caminho de escrita na primeira. Cabe registrar uma evolução em relação ao previsto. A ausência do bit de sujeira começou como ausência apenas funcional, com os registradores instanciados e constantes em zero. O parâmetro que condiciona a instanciação foi introduzido depois, e hoje a ausência é estrutural, com uma suíte de teste dedicada a verificá-la. O objetivo está alcançado, e em termos mais estritos do que o inicialmente proposto.

O quarto objetivo pedia o controlador de faltas, responsável por suspender o pipeline, sinalizar a espera ao estágio consumidor e retomar a execução. A seção~\ref{sec:5.5} descreve a máquina de estados, e a seção~\ref{sec:5.7} a integração. O ponto que merece registro é a inversão de premissa que a integração exigiu: o pipeline não consome o sinal de suspensão da cache, e deriva a espera do sinal de acerto, porque é este que marca o ciclo em que o dado está disponível. A decisão está justificada e verificada. O objetivo está alcançado.

O quinto objetivo tratava do mecanismo de consistência exigido pela instrução FENCE.I. A seção~\ref{sec:5.8} descreve as duas máquinas que o implementam, uma no pipeline e outra no coordenador, e a seção~\ref{sec:6.3} apresenta a verificação. O teste de código automodificável é o que dá sustentação ao resultado: a sub-rotina é reescrita, a sincronização é executada e a segunda chamada devolve o valor novo. Como esse desfecho exige tanto a drenagem das linhas sujas quanto a invalidação da cache de instruções, qualquer das duas omissões faria o teste falhar. O objetivo está alcançado.

O sexto objetivo pedia a verificação por testes automatizados sobre um modelo de referência. O \refCap{6_Verificacao} reporta trinta suítes e cento e oitenta casos, dos quais cento e setenta e nove aprovados, nenhum reprovado e um cancelado por ausência de variável de ambiente. Mais relevante que o número é o desenho: os casos de varredura comparam não apenas os dados, mas os contadores, contra um modelo escrito independentemente em Scala. O objetivo está alcançado, com a ressalva, registrada na seção~\ref{sec:6.8}, de que a independência disponível é entre representações, e não entre agentes.

O sétimo objetivo pedia quantificar o efeito da geometria sobre a taxa de acerto e sobre o tráfego de memória em bytes, além de determinar o custo em ciclos de cada evento do subsistema. As três partes estão cumpridas. As duas primeiras, pela varredura de treze configurações sobre um trace fixo. A terceira, pela instrumentação do custo da instrução de sincronização, que resultou em vinte medições exatas na cache isolada e dezesseis no pipeline e que, além de fornecer o número, verificou a expressão de custo da seção~\ref{sec:4.8}, corrigindo-lhe o termo fixo em quatro ciclos. O objetivo está integralmente alcançado.

\section{Síntese dos resultados}\label{sec:7.2}

A configuração adotada obteve taxa de acerto de 69,5\% sobre o trace de avaliação, com 457 faltas e 16 despejos em mil e quinhentos acessos. Isoladamente, esse número não significa nada; o que o qualifica é a varredura que o cerca.

Nos eixos de capacidade e de associatividade, os incrementos de desempenho se anulam exatamente na configuração escolhida. Elevar a capacidade de 1 KiB para 8 KiB produz ganhos sucessivos de 6,3, 7,2 e 4,1 pontos percentuais; dobrar para 16 KiB acrescenta 0,1 ponto, ou um único acerto em mil e quinhentos acessos. O mesmo padrão se repete na associatividade, com incrementos de 1,7, 0,6 e 0,1 ponto ao passar de uma para duas, quatro e oito vias. O resultado tem duas leituras. A primeira é que a geometria adotada não é arbitrária: ela está no ponto de saturação. A segunda é que o efeito mais informativo não está na taxa de acerto, e sim nos despejos, que caem de 261 para 16 ao longo do eixo de capacidade e de 65 para 8 ao longo do de associatividade. A cache de base já captura quase todas as linhas que seriam despejadas sujas, e é esse fato, mais que o acerto, que sustenta a escolha diante do custo em área.

O eixo de tamanho de linha é o único em que a configuração adotada não é a de melhor taxa de acerto, e por isso o mais informativo. A linha de 64 bytes acerta 84,2\% contra 69,5\% da de 16 bytes, uma diferença de 14,7 pontos percentuais e 220 faltas. Lido apenas por esse critério, o resultado condenaria a escolha. Lido em bytes transferidos, inverte: a linha de 64 move 19.200 bytes contra 7.568 da de 16, ou 2,5 vezes mais, para evitar 220 faltas. A conclusão que se extrai não é que uma das duas esteja errada, mas que a taxa de acerto é um critério insuficiente quando o custo da transferência é proporcional ao seu volume, e que a escolha de 16 bytes se justifica por tráfego e por área, não por acerto.

O eixo de política produziu o resultado de menor amplitude. A ordenação prevista se confirma, LRU à frente da escolha aleatória, que fica à frente da FIFO, mas a diferença entre a primeira e a última é de oito acertos em mil e quinhentos, ou 0,5 ponto percentual. Uma margem desse tamanho não distingue as três políticas de forma estatisticamente defensável, e o trabalho registra isso como achado, e não como confirmação. A justificativa da LRU repousa, ao fim, sobre a previsibilidade e sobre a possibilidade de verificá-la contra um modelo independente, que é o que permitiu tratá-la como referência na varredura.

\section{Limitações}\label{sec:7.3}

Cinco limites delimitam o alcance desses resultados, e todos foram declarados ao longo do texto.

O primeiro é a natureza da carga. O trace é sintético e construído para exibir localidade: rajadas sequenciais de uma a quatro palavras, sessenta por cento concentrados em quatro kibibytes, vinte e cinco por cento de escritas. Os números valem para esse padrão. Um trace com predomínio de saltos inverteria a conclusão do eixo de linha, porque linhas maiores passariam a trazer dados que nunca seriam usados. O que o capítulo apresenta é, portanto, uma curva de sensibilidade a parâmetros sob carga declarada, e não uma previsão de desempenho.

O segundo é o objeto medido e o modelo de memória. Cada configuração da varredura é uma cache de dados isolada, sobre memória de latência zero, e o custo de cada transferência é fixo, independentemente do seu tamanho. É essa premissa que faz da linha de 16 bytes uma escolha defensável: em um modelo que cobrasse por byte transferido, a linha maior seria penalizada; em um que cobra por transação, seria favorecida. A conclusão do trabalho está condicionada a essa premissa, e o texto o diz explicitamente.

O terceiro é a amostra. Mil e quinhentos acessos bastam para curvas com diferenças de vários pontos percentuais, e não bastam para os eixos em que a diferença é de poucos acertos. Nesses casos o resultado indica ordem de grandeza, não distinção.

O quarto é o escopo das medições. Não há comparação entre execução com e sem cache, porque o pipeline não dispõe de configuração de contorno, de modo que o trabalho não reporta ganho de desempenho em ciclos, e sim o custo de cada evento e o efeito da geometria sobre o acerto e o tráfego. Não há síntese lógica, de modo que área e frequência máxima não são quantificadas.

O quinto é a limitação de método já discutida na seção~\ref{sec:2.6}: a independência entre projeto e verificação foi obtida entre representações, e não entre agentes.

\section{Trabalhos futuros}\label{sec:7.4}

A continuidade mais direta deste trabalho está em duas etapas de custo baixo que fecham lacunas declaradas. A primeira é um modo de contorno das caches no pipeline, que tornaria possível a comparação entre execução com e sem cache e converteria o custo por evento, já medido, em ganho agregado de desempenho. A segunda é a síntese lógica, que daria aos parâmetros de área e de frequência o mesmo tratamento quantitativo que a taxa de acerto e o custo da sincronização receberam aqui.

Uma segunda frente decorre do que a análise de tráfego revelou. Se o custo da transferência for proporcional ao volume, e não ao número de transações, a linha de 16 bytes é a escolha correta e a de 64 não se justifica. Se o modelo de memória for alterado para refletir largura de banda, o compromisso se desloca, e vale reexaminá-lo. No mesmo sentido, a substituição da memória de leitura combinacional por uma de leitura síncrona reduziria o consumo de memória distribuída em arranjo de portas programáveis, ao preço de uma latência adicional a acomodar na máquina de estados.

Uma terceira frente diz respeito ao próprio controlador. A cache implementada é bloqueante, e a seção~\ref{sec:3.2} registrou que o estado da arte em cache de dados de código aberto opera com múltiplas faltas em voo. A implementação de uma cache não bloqueante, com registro de faltas em andamento, é o passo natural de evolução, e o trabalho deixa explícito por que não o deu: o ganho é pequeno em um núcleo de ordem única e sem execução especulativa, e a complexidade de verificação cresce de forma desproporcional ao escopo de uma monografia.

Por fim, a integração com o restante do núcleo condiciona o que o subsistema pode demonstrar. Enquanto não houver unidade de gerência de memória com tradução Sv32, interrupções de temporizador e entrada e saída mapeada em memória, o processador não executa um sistema operacional, e a avaliação fica restrita a programas sintéticos. A extensão da coerência para múltiplos núcleos, por sua vez, exigiria que cada um executasse a instrução de sincronização após receber interrupção do núcleo que gravou o código, uma vez que a drenagem e a invalidação locais não alcançam as caches dos demais. São essas as etapas que separam o subsistema descrito aqui do objetivo coletivo que o motivou.
